\documentclass[10pt,a4paper]{amsart}
\usepackage[margin=19mm]{geometry}
\usepackage{amsmath,amssymb,mathtools}
\usepackage[scr=rsfs,cal=euler]{mathalfa}
\usepackage{xfrac}
\usepackage[unicode,hidelinks,bookmarks=false]{hyperref}
\newcommand{\ie}{i.e.\ }
\newcommand{\cf}{cf.\ }
\newcommand{\resp}{respectively\ }
\newcommand{\Subsection}{\subsection}
\providecommand{\nobreakdash}{\nobreak\hbox{-}}
\setlength{\emergencystretch}{2em}
\multlinegap0pt
\usepackage{amssymb}
\usepackage{mathtools}
\usepackage{tikz}
\newtheorem{corollary}{Corollary}
\title{\textbf{On combinatorial bounds for the total Tjurina numbers of certain curves and surfaces with isolated singularities}}
\author{Piotr Pokora}
\date{Source: arXiv:2602.21982v5 (CC BY 4.0)}
\begin{document}
\maketitle

\noindent\emph{Source and licence.} \textbf{On combinatorial bounds for the total Tjurina numbers of certain curves and surfaces with isolated singularities}. Version arXiv:2602.21982v5 (CC BY 4.0), distributed by arXiv under the Creative Commons Attribution 4.0 International licence, http://creativecommons.org/licenses/by/4.0/, as recorded in the arXiv licence metadata for this exact version and shown on the abstract page https://arxiv.org/abs/2602.21982v5. Full licence notice: \texttt{licenses/arxiv-2602.21982-CC-BY-4.0.txt}.\par\smallskip

\noindent\textit{Editorial scope.} Counted unit: the article's corollary bbb (source0.tex lines 182--186). The article's complete original proof of it is delivered with it (source0.tex lines 187--193, one \texttt{proof} environment of 7 lines). No mathematics is written, repaired or re-derived: the statement and the proof are the article's own lines, in the article's own notation, and no retained line is substituted or re-flowed.  No further source-local context is needed: every cross-reference of every retained line resolves inside the two delivered spans.  Nothing was omitted from any retained span, so the package carries no omission record.  No retained line contains a citation, so the wrapper carries no bibliography and no bibliography entry.  No retained line contains a figure environment, an image-inclusion command or drawing code, so no image is delivered and the package declares no asset dependency.  Editorial formatting only, in addition: an AMS \texttt{amsart} preamble that restates the article's own declarations and macros used by the retained lines, namely the environments \texttt{corollary} on separate counters, together with the standard packages those lines need.  The retained environments print in this document as Corollary 1 in order of appearance, whereas the article prints the same items with the numbering of its own full document.  The complete native source of the article as distributed in the arXiv e-print for this exact version is bundled unchanged as \texttt{sources/195301/source0.tex}.
\begin{corollary}
\label{bbb}
Let $\mathcal{L} \subset \mathbb{P}^{2}$ be an arrangement of $d\geq 6$ lines such that $m(\mathcal{L})=4$. Then
$$n_{3}+3n_{4} \leq \frac{4d(d-3)}{15}.$$
\end{corollary} 

\begin{proof}
We start with the na\"ive combinatorial count:
\[ \frac{d(d-1)}{2} = n_{2} + n_{3} + n_{4} + 2n_{3}+5n_{4} \geq n_{2}+n_{3}+n_{4} + \frac{5}{3}(n_{3}+3n_{4})\geq \frac{d(d+15)}{18}+\frac{5}{3}(n_{3}+3n_{4}),\]
hence we get
$$\frac{5}{3}(n_{3}+3n_{4}) \leq \frac{8d^{2}-24d}{18},$$
and this completes the proof.
\end{proof}

\end{document}
